\documentclass[12pt,a4paper]{article}
% Drake_preamble.tex - LaTeX Preamble Template
% 作者: 謝慕揚, MD, PhD, FESC
% 
% 使用說明:
% 1. 表格請使用 longtable，不要有垂直分隔線 |
% 2. 表格內換行使用 \newline，不要使用 \\
% 3. 藥物名稱、檢驗、檢查請維持英文
% 4. Reference 格式請使用 \href 加上驗證過的 DOI

% Including essential packages for Chinese typesetting
\usepackage{xeCJK}
\usepackage{fontspec}
% Configuring Chinese font with Noto Serif CJK TC, fallback to SimSun
\IfFontExistsTF{Noto Serif CJK TC}{
  \setCJKmainfont{Noto Serif CJK TC}
  \setCJKsansfont{Noto Serif CJK TC}
  \setCJKmonofont{Noto Serif CJK TC}
}{\setCJKmainfont{SimSun}
  \setCJKsansfont{SimSun}
  \setCJKmonofont{SimSun}
}

% Including other necessary packages
\usepackage{geometry}
\usepackage{titlesec}
\usepackage{enumitem}
\usepackage{xcolor}
\usepackage{colortbl}
\usepackage{tcolorbox}
\usepackage{booktabs}
\usepackage{longtable}
\usepackage{array}
\usepackage{graphicx}
\usepackage{tikz}
\usepackage{fancyhdr}
\usepackage{multicol}
\usepackage{datetime2}
\usepackage{hyperref}
\usepackage{mdframed}
\usepackage{amsmath}
\usepackage{amssymb}
\usepackage{multirow}

\hypersetup{
    colorlinks=true,
    linkcolor=emergency_blue,
    citecolor=emergency_blue,
    urlcolor=emergency_blue,
    pdfborder={0 0 0}
}

% Defining page geometry
\geometry{
  top=2.5cm,
  bottom=2.5cm,
  left=2cm,
  right=2cm,
  headheight=25pt
}

% Adjusting topmargin to compensate for increased headheight
\addtolength{\topmargin}{-10pt}

% Defining colors
\definecolor{emergency_red}{RGB}{186,24,27}
\definecolor{emergency_blue}{RGB}{0,114,188}
\definecolor{emergency_gray}{RGB}{120,120,120}
\definecolor{highlight_yellow}{RGB}{255,250,205}

% Setting title formats
\titleformat{\section}[block]{\Large\bfseries\color{emergency_red}}{\thesection}{10pt}{}
\titleformat{\subsection}[block]{\large\bfseries\color{emergency_blue}}{\thesubsection}{8pt}{}
\titleformat{\subsubsection}[block]{\normalsize\bfseries\color{emergency_gray}}{\thesubsubsection}{6pt}{}

% Defining custom environment for important notes
\newmdenv[
    linecolor=emergency_red,
    backgroundcolor=highlight_yellow,
    linewidth=2pt,
    topline=true,
    bottomline=true,
    leftline=true,
    rightline=true,
    innertopmargin=10pt,
    innerbottommargin=10pt,
    innerleftmargin=10pt,
    innerrightmargin=10pt
]{importantbox}

% Defining note command with tcolorbox
\newcommand{\note}[1]{
    \par
    \noindent
    \begin{tcolorbox}[
        colback=emergency_blue!5!white,
        colframe=emergency_blue,
        title=\textbf{重點提醒},
        fonttitle=\bfseries,
        boxrule=0.5pt,
        arc=3pt,
        left=6pt,
        right=6pt,
        top=6pt,
        bottom=6pt
    ]
    #1
    \end{tcolorbox}
}

% Key findings box
\newcommand{\keyfinding}[1]{
    \par
    \noindent
    \begin{tcolorbox}[
        colback=yellow!10!white,
        colframe=orange!75!black,
        title=\textbf{核心發現摘要},
        fonttitle=\bfseries,
        boxrule=0.5pt,
        arc=3pt
    ]
    #1
    \end{tcolorbox}
}

% Defining custom date format for ISO 8601
\DTMnewdatestyle{iso}{
  \renewcommand{\DTMdisplaydate}[4]{\number##1-\number##2-\number##3}
  \renewcommand{\DTMDisplaydate}[4]{\number##1-\number##2-\number##3}
}
\DTMsetdatestyle{iso}
\newcommand{\mydate}{\DTMtoday}


\pagestyle{fancy}
\fancyhf{}
\fancyhead[L]{\textcolor{emergency_red}{Circulation on the Run 教學講義}}
\fancyhead[R]{\textcolor{emergency_blue}{2026年1月13日期刊重點}}
\fancyfoot[C]{\textcolor{emergency_gray}{\thepage}}
\renewcommand{\headrulewidth}{1pt}
\renewcommand{\headrule}{\hbox to\headwidth{\color{emergency_red}\leaders\hrule height \headrulewidth\hfill}}

\begin{document}

\begin{center}
{\LARGE\bfseries\textcolor{emergency_red}{Circulation on the Run Podcast}}\\[0.5cm]
{\Large\textbf{January 13, 2026 Issue Highlights}}\\[0.3cm]
{\large 2026年1月13日期刊重點整理}\\[0.5cm]
{\normalsize \href{https://doi.org/10.1161/podcast.20260112.347892}{Circulation. 2026. DOI: 10.1161/podcast.20260112.347892}}\\[0.3cm]
{\normalsize 整理：謝慕揚 MD, PhD, FESC \quad 日期：\mydate}
\end{center}

\vspace{0.5cm}

\begin{tcolorbox}[colback=yellow!10!white,colframe=orange!75!black,title=\textbf{本期核心重點}]
\begin{enumerate}
    \item \textbf{創新醫材：}中國多中心RCT證實新型\textcolor{blue}{生物可吸收PFO封堵器}在6個月及24個月追蹤時，封堵成功率與傳統nitinol device不劣性
    \item \textbf{中風二級預防：}FOURIER-OLE研究顯示，曾患缺血性中風者積極降低LDL-C至$<$40 mg/dL可降低MACE風險，且不增加出血性中風
    \item \textbf{心腎交互作用：}CKD病患的extracellular vesicles (EVs)攜帶腎臟來源microRNAs，直接損傷心肌細胞
\end{enumerate}
\end{tcolorbox}

\tableofcontents
\newpage

%============================================================
\section{專題討論：生物可吸收PFO封堵器}
%============================================================

\subsection{研究背景}

\begin{itemize}
    \item Patent foramen ovale (PFO)封堵術主要適應症為\textbf{cryptogenic stroke}預防
    \item 傳統nitinol device的問題：
    \begin{itemize}
        \item 長期留置可能誘發\textbf{心律不整}
        \item 影響後續\textbf{transseptal access}（如AF ablation、MitraClip）
        \item 年輕患者需終身攜帶金屬植入物
    \end{itemize}
    \item \textcolor{blue}{\textbf{生物可吸收裝置}}被視為此領域的「聖杯」(holy grail)
\end{itemize}

\subsection{研究設計}

\begin{longtable}{p{4cm}p{10cm}}
\toprule
\textbf{項目} & \textbf{內容} \\
\midrule
\endfirsthead
\toprule
\textbf{項目} & \textbf{內容} \\
\midrule
\endhead
\bottomrule
\endfoot

研究類型 & Prospective, multicenter, randomized controlled trial \\
\midrule
研究地點 & 中國7家醫院 \\
\midrule
收案人數 & 190位患者（試驗組96位，對照組94位） \\
\midrule
收案期間 & 2021年2月至7月 \\
\midrule
納入條件 & • 年齡18-65歲\newline
• PFO合併至少一項適應症\newline
• Cryptogenic stroke或TIA合併grade 2-3 right-to-left shunt\newline
• 儘管使用抗血小板/抗凝治療仍復發事件\newline
• Stroke/TIA合併confirmed venous thrombosis或PE \\
\midrule
主要終點 & 6個月時closure success rate（以contrast echo評估） \\
\midrule
裝置材質 & PDO monofilament + PLLA（最終代謝為CO$_2$和H$_2$O） \\
\midrule
手術導引 & 試驗組：TTE guidance（無需透視）\newline
對照組：Fluoroscopy \\
\bottomrule
\end{longtable}

\subsection{主要結果}

\begin{longtable}{p{5cm}ccc}
\toprule
\textbf{終點} & \textbf{生物可吸收組} & \textbf{Nitinol組} & \textbf{P值} \\
\midrule
\endfirsthead
\toprule
\textbf{終點} & \textbf{生物可吸收組} & \textbf{Nitinol組} & \textbf{P值} \\
\midrule
\endhead
\bottomrule
\endfoot

6個月closure success & 90\% & 91\% & NS \\
24個月closure success & 94\% & 91\% & NS \\
死亡 & 0 & 0 & -- \\
出血事件 & 0 & 0 & -- \\
心律不整 & 低且相似 & 低且相似 & NS \\
手術併發症 & 1例需外科移除 & -- & -- \\
\bottomrule
\end{longtable}

\begin{itemize}
    \item \textbf{Non-inferiority達成：}ITT分析中，兩組差異為-0.86\%，95\% CI下界高於預設margin -10\%
    \item \textbf{裝置吸收：}24個月時，TTE上hyperechoic region完全消失
    \item \textbf{無透視手術：}試驗組全程以TTE guidance完成，具發展中國家推廣潛力
\end{itemize}

\note{
\textbf{臨床意義：}此研究為生物可吸收心臟封堵器開創新局。過去類似裝置曾遇到血栓形成、不完全吸收、不對稱吸收致心律不整等問題。本研究證實可達成與傳統裝置相當的效果與安全性，且24個月後完全吸收。未來可能擴展至LAA occluder、PDA occluder等應用。
}

%============================================================
\section{臨床研究：缺血性中風後的LDL-C控制}
%============================================================

\subsection{研究背景}

\begin{itemize}
    \item MI後積極降低LDL-C的效益已確立，但\textbf{ischemic stroke}後的最佳目標值仍不明確
    \item 主要顧慮：極低LDL-C是否增加\textbf{hemorrhagic stroke}風險？
\end{itemize}

\subsection{FOURIER / FOURIER-OLE研究設計}

\begin{longtable}{p{4cm}p{10cm}}
\toprule
\textbf{項目} & \textbf{內容} \\
\midrule
\endfirsthead
\bottomrule
\endfoot

原始研究 & FOURIER：Evolocumab vs placebo於穩定ASCVD患者 \\
\midrule
追蹤期間 & FOURIER中位數2年 + FOURIER-OLE額外5年 \\
\midrule
本分析族群 & $>$5,000位曾有ischemic stroke病史的患者 \\
\midrule
主要終點 & CV death、MI、stroke、UA住院、coronary revascularization \\
\midrule
次要終點 & Stroke相關終點（包含hemorrhagic stroke） \\
\midrule
分析方法 & 依達成之LDL-C quintile分析長期結果 \\
\bottomrule
\end{longtable}

\subsection{主要發現}

\keyfinding{
在曾患缺血性中風的患者中，LDL-C降至$<$40 mg/dL可降低MACE風險（包含recurrent stroke），且\textbf{未增加hemorrhagic stroke風險}。
}

\begin{itemize}
    \item LDL-C越低，MACE風險越低，呈線性關係
    \item 達成LDL-C $<$40 mg/dL可顯著降低：
    \begin{itemize}
        \item 復發性中風
        \item 全因死亡
        \item 重大心血管事件
    \end{itemize}
    \item 重要安全性發現：\textcolor{red}{\textbf{無出血性中風風險增加的明確證據}}
\end{itemize}

\note{
\textbf{臨床建議：}此研究支持在缺血性中風二級預防時，採用更積極的LDL-C降低策略，目標可考慮$<$40 mg/dL，與冠心病二級預防策略一致。
}

%============================================================
\section{基礎研究：CKD的Extracellular Vesicles與心臟損傷}
%============================================================

\subsection{研究背景}

\begin{itemize}
    \item CKD患者心衰發生率極高，機轉不完全清楚
    \item \textbf{Extracellular vesicles (EVs)}：細胞釋放的微小顆粒，攜帶proteins、lipids、genetic material，可影響遠端器官
\end{itemize}

\subsection{研究方法與發現}

\begin{enumerate}
    \item 自CKD患者及CKD小鼠模型萃取EVs
    \item 測試EVs對心肌細胞及心臟功能的影響
    \item \textbf{主要發現：}
    \begin{itemize}
        \item CKD患者的EVs對心肌細胞有害
        \item 減少心肌收縮力、增加細胞死亡、損害心臟功能
        \item 移除CKD小鼠的EVs可改善心功能、減少心衰
        \item 有害EVs攜帶\textbf{主要來自腎臟的microRNAs}
    \end{itemize}
\end{enumerate}

\subsection{臨床意義}

\begin{tcolorbox}[colback=emergency_blue!5!white,colframe=emergency_blue,title=\textbf{心腎交互作用新機轉}]
\begin{itemize}
    \item 腎臟疾病可透過EVs直接傷害心臟
    \item EVs可作為心臟損傷的\textbf{早期預警訊號}
    \item EVs可能成為未來\textbf{治療標靶}
    \item 解釋了為何CKD患者特別容易發生HFpEF
\end{itemize}
\end{tcolorbox}

%============================================================
\section{其他本期重點文章}
%============================================================

\subsection{In-Depth Review}

\textbf{Multi-Valve Involvement in Aortic Stenosis}（Rodés-Cabau教授）
\begin{itemize}
    \item 主動脈瓣狹窄常合併其他瓣膜病變
    \item 需整體評估multiple valve disease
\end{itemize}

\subsection{Research Letter}

\textbf{Immune Checkpoint Inhibitors對心肌細胞的影響}（Wu教授）
\begin{itemize}
    \item 評估ICI在cell-derived cardiomyocytes、organoids及小鼠模型的反應
    \item 探討ICI相關心肌炎的機轉
\end{itemize}

\subsection{Letters Exchange}

\textbf{SOUL Trial相關討論}
\begin{itemize}
    \item 主題：Oral semaglutide與cardiovascular outcomes
    \item 依SGLT-2 inhibitor使用狀況的pre-specified分析
\end{itemize}

\subsection{Path to Discovery}

\textbf{致敬Francois Abboud教授}
\begin{itemize}
    \item 心血管神經與發炎調控領域的先驅科學家
    \item 數十年來培育無數後進，影響深遠
\end{itemize}

%============================================================
\section{縮寫對照表}
%============================================================

\begin{longtable}{p{4cm}p{10cm}}
\toprule
\textbf{縮寫} & \textbf{全名} \\
\midrule
\endfirsthead
\bottomrule
\endfoot

PFO & Patent Foramen Ovale (卵圓孔未閉) \\
TTE & Transthoracic Echocardiography (經胸心臟超音波) \\
PLLA & Poly-L-Lactic Acid (聚左旋乳酸) \\
PDO & Polydioxanone (聚對二氧環己酮) \\
LDL-C & Low-Density Lipoprotein Cholesterol (低密度脂蛋白膽固醇) \\
MACE & Major Adverse Cardiovascular Events (重大心血管不良事件) \\
EVs & Extracellular Vesicles (細胞外囊泡) \\
CKD & Chronic Kidney Disease (慢性腎臟病) \\
HFpEF & Heart Failure with Preserved Ejection Fraction (射出分率保留型心衰) \\
ICI & Immune Checkpoint Inhibitors (免疫檢查點抑制劑) \\
ASCVD & Atherosclerotic Cardiovascular Disease (動脈粥狀硬化性心血管疾病) \\
\bottomrule
\end{longtable}

%============================================================
\section{參考文獻}
%============================================================

\begin{enumerate}
    \item Pan X, Dong J, et al. Transcatheter Closure of Patent Foramen Ovale With a Novel Biodegradable Device: A Prospective, Multicenter, Randomized Controlled Clinical Trial. \href{https://www.ahajournals.org/doi/10.1161/CIRCULATIONAHA.125.074609}{\textit{Circulation}. 2026.}

    \item Giugliano RP, et al. Long-term LDL-C lowering in patients with prior ischemic stroke: FOURIER and FOURIER-OLE. \textit{Circulation}. 2026.

    \item Sahoo S, et al. Extracellular vesicles from chronic kidney disease patients impair cardiac function. \textit{Circulation}. 2026.

    \item Rodés-Cabau J, et al. Multi-Valve Involvement in Aortic Stenosis: Insights from a Narrative Review. \textit{Circulation}. 2026.
\end{enumerate}

\end{document}
